\chapter{Ejemplos de investigación: cómo reaparece la teoría en Haseman y operadores integrales singulares}

\section{De la condición de frontera al operador}

En Karlovich \cite[pp.~463--464, ecs.~(1.1)--(1.2)]{Karlovich2017}, $\Gamma$ es una
estrella abierta de espirales logarítmicas y $1<p<\infty$.
Se toman $G\in\SO(\Gamma)$, $g\in L^p(\Gamma)$ y un difeomorfismo
$\alpha:\Gamma\to\Gamma$ que preserva la orientación de cada rama y
satisface las hipótesis de la fuente. La condición de Haseman es
\[
 \Phi^+(\alpha(t))=G(t)\Phi^-(t)+g(t),\qquad t\in\Gamma.
\]
Aquí $S_\Gamma\in\mathcal B(L^p(\Gamma))$ es el operador integral singular de Cauchy,
$V_\alpha f=f\circ\alpha$ y $M_Gf=Gf$. Si $\varphi$ es la densidad de Cauchy, las
fórmulas de Sokhotski--Plemelj dan
\[
 \Phi^+=P_+\varphi,\qquad \Phi^-=-P_-\varphi,\qquad
 P_\pm=\tfrac12(I\pm S_\Gamma).
\]
El signo de la traza negativa explica el signo más en la ecuación operatorial:
\[
 \bigl(V_\alpha P_++M_GP_-\bigr)\varphi=g,\qquad
 T:=V_\alpha P_++M_GP_-\in\mathcal B(L^p(\Gamma)).
\]
El curso proporciona el lenguaje para preguntar si $T$ es invertible y estudiar su
espectro. La reducción de frontera y el criterio de Fredholm son resultados posteriores
que dependen de las hipótesis geométricas y de símbolos del artículo.

\begin{cthejemplo}{(La identidad como caso de la ecuación de Haseman)}
En $X=L^p(\Gamma)$ consideremos la fórmula $T=V_\alpha P_++M_GP_-$
obtenida arriba. Queremos comprobar qué operador resulta cuando
$\alpha=\mathrm{id}_\Gamma$ y $G=1$. En ese caso $V_\alpha=I_X$ y
$M_G=I_X$, mientras que la definición de $P_\pm$ da
\[
 P_++P_- =\tfrac12(I_X+S_\Gamma)+\tfrac12(I_X-S_\Gamma)=I_X.
\]
Por tanto $T=P_++P_-=I_X$. Este cálculo docente muestra que la
reducción por Plemelj se puede leer sin suponer $S_\Gamma^2=I_X$.
\end{cthejemplo}

\begin{center}
\begin{tikzpicture}[node distance=5mm and 8mm]
\node[theorybox,text width=3.35cm] (b) {Frontera de Haseman\\$\Phi^+\circ\alpha=G\Phi^-+g$};
\node[theorybox,text width=3.35cm,right=of b] (p) {Sokhotski--Plemelj\\$\Phi^+=P_+\varphi$, $\Phi^-=-P_-\varphi$};
\node[researchbox,text width=3.35cm,right=of p] (t) {$T=V_\alpha P_++M_GP_-$\\$T\varphi=g$};
\draw[arr] (b)--(p);
\draw[arr] (p)--(t);
\end{tikzpicture}
\end{center}

\section{Del ideal compacto a la propiedad de Fredholm}

Sea $X$ un espacio de Banach complejo infinito-dimensional. El ideal bilateral cerrado
$\mathcal K(X)\triangleleft\mathcal B(X)$ produce la álgebra de Calkin
$\mathcal Q(X)=\mathcal B(X)/\mathcal K(X)$. Para $A\in\mathcal B(X)$,
\[
 A\text{ Fredholm}
 \quad\Longleftrightarrow\quad
 A+\mathcal K(X)\in\GL(\mathcal Q(X)).
\]
Karlovich enuncia esta equivalencia en la introducción de su artículo de 2025
\cite[p.~2]{Karlovich2025}. La Tarea 1.4 explica por qué el cociente por un ideal
cerrado sigue siendo una álgebra de Banach; la equivalencia de Fredholm añade un
teorema propio de operadores.

\begin{cthejemplo}{(Invertible módulo compactos sin inversa exacta)}
El criterio de Atkinson anterior distingue la inversa exacta de una
inversa en Calkin. Para verlo en $X=\ell^2$, tome $P$ la proyección ortogonal de rango uno sobre
$\mathbb Ce_1$ y $A=I_X-P$. Se tiene
$\ker A=\operatorname{Ran}P\ne\{0\}$, así que $A$ no es invertible. Sin embargo,
$P\in\mathcal K(X)$ y $A+\mathcal K(X)=I_X+\mathcal K(X)$: su clase es la
unidad de $\mathcal Q(X)$. Además,
$\operatorname{Ran}(A)=(\mathbb Ce_1)^\perp$ es cerrado, mientras que el núcleo
y el conúcleo tienen dimensión uno. Así $A$ es Fredholm. El cálculo es una
elaboración docente del criterio anterior.
\end{cthejemplo}

Si se define el \emph{espectro esencial de Fredholm} por
\[
 \sigma_{\mathrm{ess}}(A)
 :=\sigma_{\mathcal Q(X)}\bigl(A+\mathcal K(X)\bigr),
\]
entonces $0\notin\sigma_{\mathrm{ess}}(A)$ equivale a que $A$ sea Fredholm.
Esto no dice que $0\notin\sigma_{\mathcal B(X)}(A)$; el ejemplo $I_X-P$
muestra la diferencia.

\begin{center}
\begin{tikzpicture}[node distance=5mm and 7mm]
\node[theorybox,text width=2.65cm] (ideal) {$\mathcal K(X)\triangleleft\mathcal B(X)$\\ideal cerrado};
\node[theorybox,text width=2.65cm,right=of ideal] (quot) {$\mathcal Q(X)$\\cociente};
\node[theorybox,text width=2.65cm,right=of quot] (class) {$A+\mathcal K(X)$\\clase invertible};
\node[researchbox,text width=2.65cm,right=of class] (fred) {$A$ Fredholm\\criterio de Calkin};
\draw[arr] (ideal)--(quot);
\draw[arr] (quot)--(class);
\draw[arr] (class)--(fred);
\end{tikzpicture}
\end{center}

En un espacio de Hilbert $H$, $\mathcal B(H)$ y su cociente por
$\mathcal K(H)$ son $C^*$-álgebras. Para un Banach general, el argumento
anterior sólo necesita la estructura de álgebra de Banach. Los cálculos de
símbolos y la localización intentan convertir la invertibilidad de la clase
en condiciones verificables para una \emph{álgebra de operadores específica};
no basta con escribir el cociente para obtener un criterio.

\Needspace{17\baselineskip}
\section{Proyecciones de Cauchy: el caso modelo y el caso Haseman}

La comparación precisa qué hipótesis convierte la fórmula
$(I\pm S_\Gamma)/2$ en una proyección y qué parte de ella se utiliza en
la ecuación de Haseman.

En el modelo $\Gamma=\mathbb R$ o $\mathbb T$ sobre $L^2(\Gamma)$,
sea $S_\Gamma$ el operador integral singular de Cauchy.
Bastos--Fernandes--Karlovich recuerdan $S_\Gamma^2=I$ y que
\[
 P_\Gamma^\pm=\tfrac12(I\pm S_\Gamma)
\]
son proyecciones ortogonales \cite[§2, p.~22]{BFK2006}. La identidad
$S_\Gamma^2=I$ da $P_\Gamma^\pm{}^2=P_\Gamma^\pm$ y
$P_\Gamma^+P_\Gamma^-=0$; la autoadjunción de $S_\Gamma$ da además
$(P_\Gamma^\pm)^*=P_\Gamma^\pm$. Es ésta la conexión literal con las
proyecciones $p=p^*=p^2$ de una $C^*$-álgebra.

\begin{center}
\begin{tikzpicture}[node distance=5mm and 7mm]
\node[theorybox,text width=2.6cm] (s) {$S_\Gamma^2=I$\\$S_\Gamma^*=S_\Gamma$};
\node[theorybox,text width=2.6cm,right=of s] (def) {$P_\Gamma^\pm=(I\pm S_\Gamma)/2$};
\node[theorybox,text width=2.6cm,right=of def] (proj) {$P_\Gamma^\pm{}^2=P_\Gamma^\pm$\\$(P_\Gamma^\pm)^*=P_\Gamma^\pm$};
\node[researchbox,text width=2.6cm,right=of proj] (split) {$L^2(\Gamma)=\operatorname{Ran}P_\Gamma^+\oplus\operatorname{Ran}P_\Gamma^-$};
\draw[arr] (s)--(def);
\draw[arr] (def)--(proj);
\draw[arr] (proj)--(split);
\end{tikzpicture}
\end{center}

En la estrella abierta de Karlovich 2017, la misma \emph{fórmula}
$P_\pm=(I\pm S_\Gamma)/2$ aparece en Sokhotski--Plemelj y en
$T=V_\alpha P_++M_GP_-$, pero no se puede importar la idempotencia del
modelo anterior. Si $X$ es un espacio de Banach, $S\in\mathcal B(X)$
y $P_\pm=(I_X\pm S)/2$, entonces
\[
 P_\pm^2-P_\pm=\tfrac14(S^2-I_X),\qquad
 P_+P_-=\tfrac14(I_X-S^2),\qquad P_++P_-=I_X.
\]
De hecho, el símbolo Mellin de $S_\Gamma$ en el caso de una sola rama,
$p=2$ y parámetro geométrico $\delta=0$ es
$\varepsilon_1\tanh(\pi x)$, donde $\varepsilon_1\in\{-1,1\}$
registra la orientación; su cuadrado
no es idénticamente uno \cite[teorema~5.3, p.~476]{Karlovich2017}.
Por tanto, en ese caso $P_\pm$ no son proyecciones algebraicas.
La conexión correcta con Haseman es la \emph{descomposición de las trazas}
por Plemelj, no una descomposición ortogonal del espacio.

\begin{cthnota}{(Dónde intervienen las hipótesis)}
La igualdad $S_\Gamma^2=I$ es exactamente la que da idempotencia;
$S_\Gamma^*=S_\Gamma$ es necesaria para ortogonalidad. El modelo
$L^2(\mathbb R)$ o $L^2(\mathbb T)$ satisface ambas propiedades en la
fuente citada. El problema de Haseman de 2017 actúa en $L^p(\Gamma)$
sobre una estrella abierta y sólo usa aquí la fórmula de Plemelj.
\end{cthnota}

\Needspace{25\baselineskip}
\section{Shifts unitarios y shifts de composición}

Comparamos una rotación que es unitaria con el operador de composición
que aparece en Haseman. El espacio de acción y la medida deciden qué
propiedad se puede afirmar.

En \cite[§1, p.~20]{BFK2006}, el círculo $\mathbb T$ actúa mediante una representación unitaria;
para $z,t\in\mathbb T$ y $\varphi\in L^2(\mathbb T)$ se escribe
\[
 U:\mathbb T\to\mathcal B(L^2(\mathbb T)),
 \qquad
 (U_z\varphi)(t)=\varphi(zt).
\]
La rotación preserva la medida de Haar $m$: al sustituir $u=zt$,
\[
 \|U_z\varphi\|_2^2=\int_{\mathbb T}|\varphi(zt)|^2\,dm(t)
 =\int_{\mathbb T}|\varphi(u)|^2\,dm(u)=\|\varphi\|_2^2.
\]
Como $U_z^{-1}=U_{z^{-1}}$, $U_z$ es unitario. En el problema de Haseman aparece en cambio
$V_\alpha f=f\circ\alpha$ sobre $L^p(\Gamma)$; Karlovich impone condiciones
de regularidad y control de la derivada del shift en cada rama
\cite[§2, p.~467]{Karlovich2017}. Estas condiciones permiten trabajar con
un operador acotado, pero no afirman que sea unitario.

\begin{center}
\begin{tikzpicture}[node distance=10mm]
\node[theorybox] (shift) {Idea común\\operador de composición / shift};
\node[researchbox,above right=of shift] (unit) {$U_z$ en $L^2(\mathbb T)$\\unitario};
\node[researchbox,below right=of shift] (hasa) {$V_\alpha$ en $L^p(\Gamma)$\\acotado bajo hipótesis del shift};
\node[theorybox,right=22mm of unit] (uspec) {Unitario $\Rightarrow\spec(U_z)\subset\mathbb T$};
\node[theorybox,right=22mm of hasa] (fred) {Se combina con $P_\pm$ y $M_G$ en el operador de Haseman};
\draw[arr] (shift)--(unit);
\draw[arr] (shift)--(hasa);
\draw[arr] (unit)--(uspec);
\draw[arr] (hasa)--(fred);
\end{tikzpicture}
\end{center}

\section{Funciones lentamente oscilantes y fibras de caracteres}

Para $a\in C_b(\mathbb R_+)$, escriba
$\operatorname{osc}(a,[r,2r])=\sup_{s,t\in[r,2r]}|a(s)-a(t)|$.
Karlovich define $\SO(\mathbb R_+)$ exigiendo que esta oscilación tienda a
cero tanto cuando $r\to0$ como cuando $r\to\infty$
\cite[§2, pp.~465--466]{Karlovich2017}. Para $I=[r,2r]$ y funciones
$a,b\in C_b(\mathbb R_+)$ se tiene
\[
\begin{aligned}
 \operatorname{osc}(a+b,I)&\le\operatorname{osc}(a,I)+\operatorname{osc}(b,I),\\
 \operatorname{osc}(ab,I)&\le\|a\|_\infty\operatorname{osc}(b,I)
                         +\|b\|_\infty\operatorname{osc}(a,I),\\
 \operatorname{osc}(\overline a,I)&=\operatorname{osc}(a,I).
\end{aligned}
\]
Si $a_n\to a$ uniformemente, entonces
$\operatorname{osc}(a,I)\le2\|a-a_n\|_\infty+
\operatorname{osc}(a_n,I)$; primero se fija $n$ grande y después se
deja tender $r$ al extremo. Así la clase es una subálgebra cerrada,
contiene las constantes y es estable bajo la involución puntual.
La identidad $\|\overline a a\|_\infty=\|a\|_\infty^2$ muestra que
es una $C^*$-álgebra conmutativa unital.

\begin{cthejemplo}{(Oscilación lenta sin límite ordinario)}
Queremos comprobar que la oscilación lenta en los extremos no obliga a
tener un límite ordinario. En $\mathbb R_+$ consideremos la función
\[
 a(r)=\sin\!\bigl(\log(1+|\log r|)\bigr),\qquad r>0,
\]
pertenece a $\SO(\mathbb R_+)$. Si $s,t\in[r,2r]$, entonces
$|\log s-\log t|\le\log2$. Ponga $h(u)=\log(1+|u|)$.
En el intervalo $I_r=[\log r,\log r+\log2]$, cuando está alejado de $0$,
el teorema del valor medio da
\[
 |h(\log s)-h(\log t)|
 \le \frac{\log2}{1+\min_{u\in I_r}|u|}.
\]
El denominador tiende a infinito tanto si $r\to0$ como si $r\to\infty$.
Como $|\sin v-\sin w|\le|v-w|$, la oscilación de $a$ en $[r,2r]$
tiende a cero en ambos extremos. Para ver que no hay límite ordinario,
defina
\[
 r_n^+=\exp\!\bigl(e^{\pi/2+2\pi n}-1\bigr),\qquad
 r_n^-=\exp\!\bigl(e^{3\pi/2+2\pi n}-1\bigr).
\]
Ambas sucesiones tienden a infinito, mientras que
$a(r_n^+)=1$ y $a(r_n^-)=-1$. Sus recíprocas tienden a cero y dan
los mismos dos valores, porque $a(r)=a(1/r)$.
Esta función, elaborada a partir de la definición de la fuente, muestra por
qué los caracteres de frontera de $\SO(\mathbb R_+)$ no se pueden sustituir
simplemente por evaluaciones $a(0)$ y $a(\infty)$.
\end{cthejemplo}

El espacio maximal $\Omega(\SO(\mathbb R_+))$ contiene las evaluaciones
en $r>0$ y fibras $M_0,M_\infty$ sobre los extremos
\cite[§2, p.~466]{Karlovich2017}. La transformada de Gelfand
$a\mapsto\widehat a$, $\widehat a(\xi)=\xi(a)$, permite describir
valores de estos caracteres sin fingir que existen $a(0)$ y
$a(\infty)$. En los cálculos de símbolos, esas fibras son parte del
espacio de parámetros; el mapa de Gelfand de los coeficientes y el
símbolo completo de un operador son mapas distintos.

\section{Generadores, cocientes y mapas de símbolos}

Karlovich \cite[ec.~(1.4), p.~3]{Karlovich2025} fija una curva estrellada
con cúspides $\Gamma$ de (1.1)--(1.2), un peso de potencia admisible $w$
de su §2.1, con $1<p<\infty$ y exponente
$-1/p<\alpha_{\mathrm{peso}}<1-1/p$, y considera la álgebra
de Banach generada por
\[
 \mathfrak B=
 \operatorname{alg}_{\mathbb R}
 \{M_a,S_\Gamma,C:\ a\in\SO^\diamond(\Gamma)\}
 \subset\mathcal B(\mathbb RL^p(\Gamma,w)).
\]
Aquí $M_a$ multiplica por $a$, $S_\Gamma$ es el operador singular de
Cauchy y $Cf=\overline f$ es conjugación compleja. El símbolo
$\mathbb RL^p(\Gamma,w)$ significa el espacio $L^p$ considerado sobre
los escalares reales: para $f\in L^p(\Gamma,w)$ se tiene
$C(if)=\overline{if}=-i\overline f=-iCf$, mientras que
$C(rf)=rCf$ para $r\in\mathbb R$. Así $C$ es lineal sobre
$\mathbb R$, pero no sobre $\mathbb C$. Por eso esta álgebra no se convierte automáticamente en una
$C^*$-álgebra compleja ni se puede aplicar sin más el teorema espectral
complejo de la Parte 1. La notación $\SO^\diamond(\Gamma)$ pertenece
al artículo de 2025 y no designa la misma álgebra
$\SO(\mathbb R_+)$ del artículo de 2017.
El artículo de 2025 llama $\alpha$ al exponente del peso; aquí el
subíndice evita confundirlo con el shift $\alpha$ de Karlovich 2017.

\begin{center}
\begin{tikzpicture}[node distance=5mm and 8mm]
\node[researchbox,text width=3.25cm] (gen) {$M_a,\ S_\Gamma,\ C$\\generadores tipados};
\node[theorybox,text width=3.25cm,right=of gen] (alg) {$\mathfrak B$\\cierre real bajo suma,\\producto y norma};
\node[researchbox,text width=3.25cm,right=of alg] (sym) {Cociente y símbolos\\criterio en una álgebra fijada};
\draw[arr] (gen)--(alg);
\draw[arr] (alg)--(sym);
\end{tikzpicture}
\end{center}

La Tarea 1.3 enseña qué significa generar y cerrar una álgebra. La
Tarea 1.4 enseña a pasar a un cociente. En el artículo de 2025, el
paso adicional es construir mapas matriciales de símbolo y probar
que detectan Fredholmness. El teorema~6.2 trata la subálgebra compleja
$\mathfrak A=\operatorname{alg}_{\mathbb C}
\{M_a,S_\Gamma,CS_\Gamma C:a\in\SO^\diamond(\Gamma)\}$;
la representación matricial compleja del teorema~7.1 permite después formular el criterio
para $N\in\mathfrak B$ mediante el mapa de símbolo $Y$ del teorema~7.2
\cite[pp.~25--26]{Karlovich2025}. Este último exige pertenencia a
$\mathfrak B$ y las hipótesis de curva, peso y coeficientes del artículo;
no certifica por sí solo un operador de Haseman con shift.

\begin{cthnota}{(Cuatro mapas relacionados y diferentes)}
Un homomorfismo de álgebras preserva suma y producto. Un
$*$-homomorfismo preserva además la involución. La transformada de
Gelfand $a\mapsto\widehat a$ de una álgebra conmutativa evalúa en sus
caracteres. Un mapa de símbolo de Fredholm se construye para una
álgebra concreta de operadores, frecuentemente después de pasar
módulo compactos, y puede tomar valores matriciales. Estas
construcciones comparten operaciones algebraicas, pero sus dominios,
codominios e hipótesis no son intercambiables.
\end{cthnota}

\Needspace{20\baselineskip}
\section{Ruta completa desde el problema hasta un criterio}

En el diagrama, azul marca el lenguaje del curso, verde el objeto de
investigación y la flecha hacia símbolos presupone que se ha construido
el cálculo correspondiente para la álgebra concreta. Aquí
$X=L^p(\Gamma)$, de modo que $T+\mathcal K(X)$ es una clase de la
álgebra de Calkin asociada al mismo espacio de acción. El criterio de
Karlovich 2017 exige además las hipótesis de su teorema 7.1
\cite[pp.~487--489]{Karlovich2017}.

\begin{center}
\begin{tikzpicture}[node distance=5mm and 7mm]
\node[researchbox,text width=3.25cm] (h) {[investigación]\\frontera de Haseman};
\node[researchbox,text width=3.25cm,right=of h] (p) {[investigación]\\Sokhotski--Plemelj};
\node[researchbox,text width=3.25cm,right=of p] (t) {[investigación]\\$T=V_\alpha P_++M_GP_-$};
\node[theorybox,text width=3.25cm,below=of t] (b) {[curso]\\$T\in\mathcal B(L^p(\Gamma))$};
\node[theorybox,text width=3.25cm,left=of b] (q) {[curso + puente]\\$T+\mathcal K(X)$ en $\mathcal Q(X)$};
\node[researchbox,text width=3.25cm,left=of q] (f) {[puente]\\invertibilidad de clase\\$\Longleftrightarrow$ Fredholm};
\node[researchbox,text width=3.25cm,below=of f] (s) {[investigación]\\símbolos y localización};
\node[researchbox,text width=3.25cm,right=of s] (c) {[investigación]\\criterio de Fredholm\\bajo hipótesis del artículo};
\draw[arr] (h)--(p);
\draw[arr] (p)--(t);
\draw[arr] (t)--(b);
\draw[arr] (b)--(q);
\draw[arr] (q)--(f);
\draw[arr] (f)--(s);
\draw[arr] (s)--(c);
\end{tikzpicture}
\end{center}

La última fila describe cómo se \emph{comprueba} Fredholmness en la
literatura. No afirma que todo operador Fredholm tenga automáticamente
el símbolo concreto del artículo, ni que el criterio de 2017 se aplique
a un operador con cúspides de 2025.

\Needspace{19\baselineskip}
\section{Una ruta de lectura para el manuscrito de investigación}

Cuando aparezca una construcción operatorial en el manuscrito, conviene hacer las siguientes
preguntas en este orden:
\begin{enumerate}
\item ¿En qué álgebra de Banach vive el operador?
\item ¿Cuál es la unidad y cuáles son los generadores relevantes?
\item ¿Se está trabajando exactamente o módulo un ideal, normalmente $\K(X)$?
\item ¿La subálgebra considerada es conmutativa o sólo posee una subálgebra central conmutativa?
\item ¿Aparecen caracteres, espacio maximal o transformada de Gelfand?
\item ¿Hay una estructura de $C^*$-álgebra compleja (típicamente en $L^2$) que permita usar normalidad, unitarios,
      proyecciones o cálculo funcional?
\item ¿Qué parte del criterio de Fredholm equivale a invertibilidad en un cociente o de un símbolo?
\end{enumerate}

Estas preguntas obligan a distinguir el lenguaje común de las hipótesis que autorizan
cada teorema. Por ejemplo, un criterio publicado para una álgebra de operadores con cúspides
no se aplica automáticamente a un operador con shift de otra geometría.
